%% Compétences : 4C04-STAT01, 4C04-STAT02, 4C04-STAT03
%% Cours (trace écrite copiée au tableau) : une partie par notion.

%% Font size %%
\documentclass[11pt]{article}

%% Load the custom package
\usepackage{Mathdoc}
\usepackage{icomma}
\usepackage{colortbl}

%% Numéro de séquence %% Titre de la séquence %%
\renewcommand{\centerhead}{S04 - Statistiques}

% Encadré « compétence » en tête de partie
\newcommand{\comp}[2]{\par\smallskip\noindent\colorbox{gray!15}{\parbox{\dimexpr\linewidth-2\fboxsep}{\small\textbf{Compétence #1} : #2}}\par\smallskip}
\newcommand{\f}[2]{\dfrac{#1}{#2}}

%% Spacing commands %%
\renewcommand{\baselinestretch}{1} \setlength{\parindent}{0pt}

\begin{document}

\begin{tcolorbox}[colback=white,colframe=black,boxrule=0.6pt,arc=2mm,title={\textbf{Ce que je dois savoir faire à la fin du chapitre}},coltitle=black,colbacktitle=gray!15]
\begin{tabular}{@{}ll@{}}
\textbf{4C04-STAT01} & Calculer l'\textbf{étendue} et une \textbf{fréquence} \hfill (parties 2 et 3)\\
\textbf{4C04-STAT02} & Calculer la \textbf{moyenne} d'une série (simple ou pondérée) \hfill (partie 4)\\
\textbf{4C04-STAT03} & Déterminer la \textbf{médiane} d'une série \hfill (partie 5)\\
\end{tabular}
\end{tcolorbox}

%% ------------------------------------------------------------------
\section{Vocabulaire}

\begin{vocabulaire}
On interroge les élèves d'une classe sur leur nombre de frères et sœurs.
\begin{itemize}
\item la \textbf{population} étudiée : les élèves de la classe ;
\item le \textbf{caractère} étudié : le nombre de frères et sœurs ;
\item les \textbf{valeurs} : $0$, $1$, $2$, $3$\ldots ;
\item l'\textbf{effectif} d'une valeur : le nombre de fois où elle apparaît ;
l'\textbf{effectif total} : le nombre de données.
\end{itemize}
\end{vocabulaire}

\begin{center}
\renewcommand{\arraystretch}{1.4}
\begin{tabular}{|l|c|c|c|c|c||c|}
\hline
\textbf{Nombre de frères et sœurs} & $0$ & $1$ & $2$ & $3$ & $4$ & \textbf{Total} \\ \hline
\textbf{Effectif} & $4$ & $9$ & $6$ & $3$ & $2$ & $24$ \\ \hline
\end{tabular}
\end{center}

%% ------------------------------------------------------------------
\section{L'étendue}
\comp{4C04-STAT01}{calculer l'étendue \hfill \textit{Fiche Ex-4C04-01}}

\begin{definition}
L'\textbf{étendue} d'une série est la différence entre la \textbf{plus grande} valeur et la \textbf{plus petite} valeur :
\qquad étendue $=$ valeur maximale $-$ valeur minimale.
\end{definition}

\begin{exemple}
Notes de Léo : $12$ ; $8$ ; $15$ ; $6$ ; $11$. \quad Étendue : $15 - 6 = 9$ points.\\
L'étendue mesure la \textbf{dispersion} : plus elle est grande, plus les valeurs sont « étalées ».
\end{exemple}

%% ------------------------------------------------------------------
\section{La fréquence}
\comp{4C04-STAT01}{calculer une fréquence \hfill \textit{Fiche Ex-4C04-02}}

\begin{definition}
La \textbf{fréquence} d'une valeur est la proportion de cette valeur dans la série :
\[ \text{fréquence} = \f{\text{effectif de la valeur}}{\text{effectif total}} \]
On la donne en \textbf{fraction}, en \textbf{nombre décimal} ou en \textbf{pourcentage}.
\end{definition}

\begin{exemple}
Dans le tableau de la partie 1, la fréquence de la valeur $2$ est
$\f{6}{24} = 0,25 = 25\,\%$ : un quart des élèves a $2$ frères et sœurs.
\end{exemple}

\begin{remarque}
Une fréquence est toujours comprise entre $0$ et $1$ (entre $0\,\%$ et $100\,\%$).
La somme de toutes les fréquences vaut $1$ (soit $100\,\%$).
\end{remarque}

%% ------------------------------------------------------------------
\section{La moyenne}
\comp{4C04-STAT02}{calculer une moyenne \hfill \textit{Fiches Ex-4C04-03 et 04}}

\begin{definition}[Moyenne simple]
\[ \text{moyenne} = \f{\text{somme de toutes les valeurs}}{\text{nombre de valeurs}} \]
\end{definition}

\begin{exemple}
Notes de Léo : $12$ ; $8$ ; $15$ ; $6$ ; $11$. \quad
Moyenne : $\f{12 + 8 + 15 + 6 + 11}{5} = \f{52}{5} = 10,4$.
\end{exemple}

\begin{definition}[Moyenne pondérée]
Quand les données sont dans un tableau d'effectifs, on multiplie chaque valeur par son effectif :
\[ \text{moyenne} = \f{\text{valeur}_1 \times \text{effectif}_1 + \text{valeur}_2 \times \text{effectif}_2 + \cdots}{\text{effectif total}} \]
\end{definition}

\begin{exemple}
Avec le tableau de la partie 1 :
$\f{0\times 4 + 1\times 9 + 2\times 6 + 3\times 3 + 4\times 2}{24} = \f{38}{24} \approx 1,6$ frère ou sœur en moyenne.
\end{exemple}

%% ------------------------------------------------------------------
\section{La médiane}
\comp{4C04-STAT03}{déterminer la médiane \hfill \textit{Fiches Ex-4C04-05 et 06}}

\begin{definition}
La \textbf{médiane} d'une série \textbf{rangée dans l'ordre croissant} partage la série en \textbf{deux groupes de même effectif} :
au moins la moitié des valeurs sont inférieures ou égales à la médiane, au moins la moitié sont supérieures ou égales.
\end{definition}

\begin{propriete}[Méthode]
\begin{enumerate}
\item Je range les valeurs dans l'ordre croissant et je compte l'effectif total $N$.
\item Si $N$ est \textbf{impair} : la médiane est la valeur du milieu, de rang $\f{N+1}{2}$.
\item Si $N$ est \textbf{pair} : la médiane est la moyenne des deux valeurs du milieu, de rangs $\f{N}{2}$ et $\f{N}{2}+1$.
\end{enumerate}
\end{propriete}

\begin{exemple}
\textbf{$N$ impair} : $6$ ; $8$ ; $\mathbf{11}$ ; $12$ ; $15$ \quad ($N = 5$, rang $3$) \quad médiane : $11$.\\
\textbf{$N$ pair} : $3$ ; $7$ ; $\mathbf{9}$ ; $\mathbf{13}$ ; $14$ ; $20$ \quad ($N = 6$, rangs $3$ et $4$) \quad médiane : $\f{9 + 13}{2} = 11$.
\end{exemple}

\begin{exemple}[Avec un tableau d'effectifs : les effectifs cumulés]
Tableau de la partie 1 : $N = 24$, la médiane est entre la $12^\text{e}$ et la $13^\text{e}$ valeur.
\begin{center}
\renewcommand{\arraystretch}{1.3}
\begin{tabular}{|l|c|c|c|c|c|}
\hline
\textbf{Valeur} & $0$ & $1$ & $2$ & $3$ & $4$ \\ \hline
\textbf{Effectif} & $4$ & $9$ & $6$ & $3$ & $2$ \\ \hline
\textbf{Effectif cumulé} & $4$ & $13$ & $19$ & $22$ & $24$ \\ \hline
\end{tabular}
\end{center}
Les valeurs de rang $5$ à $13$ valent $1$ : la $12^\text{e}$ et la $13^\text{e}$ valent $1$, donc la médiane est $1$.
\end{exemple}

\begin{remarque}
\textbf{Moyenne ou médiane ?} La moyenne est sensible aux valeurs extrêmes, pas la médiane.
Salaires dans une entreprise : $1\,400$ ; $1\,500$ ; $1\,600$ ; $1\,700$ ; $9\,000$ €.
Moyenne : $3\,040$ € ; médiane : $1\,600$ €, plus représentative de la plupart des salariés.
\end{remarque}

\end{document}

%%% Local Variables:
%%% mode: LaTeX
%%% TeX-master: t
%%% End:
